\documentclass[../main.tex]{subfiles}

\begin{document}

\chapter{ Fundemental Group }

\begin{definition}{}{}
Let $X$ be a topological space. Let $\omega, \tau: \mathbb{I} \longrightarrow X$ be two paths such that the terminal point $\omega(1)$ of $\omega$ is the same as the initial point $\tau(0)$ of $\tau$. We then define $\omega * \tau: \mathbb{I} \longrightarrow X$ by the formula
\begin{equation}
\omega * \tau(t)=\left\{\begin{array}{ll}
\omega(2 t), & 0 \leq t \leq 1 / 2, \\
\tau(2 t-1), & 1 / 2 \leq t \leq 1 .
\end{array}\right.
\end{equation}
\end{definition}
\begin{lemma}{}{}
If $\omega_{1} \simeq \omega_{2}$ and $\tau_{1} \simeq \tau_{2}$ and $\omega_{1} * \tau_{1}$ is defined then $\omega_{2} * \tau_{2}$ is defined and we have, $\omega_{1} * \tau_{1} \simeq \omega_{2} * \tau_{2}$.
\end{lemma}

\begin{definition}{}{}
    We can use the standard parameterisation of any closed interval $[a, b]$ by the interval $[0, 1]$ and think of a map $\gamma : [a, b] \to X$ as a path in $X$, viz.,
\[
\hat{\gamma}(t) := \gamma(((b - a)t + a)).
\]
\end{definition}

\begin{theorem}{Invariance under subdivision}{}
Let $\gamma : [0,1] \to X$ be a path, and $0 < t_1 < \dots < t_n < 1$. Then $\gamma$ is path homotopic to $\widehat{\gamma|_{[0,t_1]}} * \dots * \widehat{\gamma|_{[t_n,1]}}$.
\end{theorem}

\begin{definition}{The Fundamental Group}{}
Let $X$ be any topological space and $x_0 \in X$. By the above two lemmas it follows that the set $\pi_1(X, x_0)$ consisting of all path-homotopy classes of loops in $X$ based at $x_0$ forms a group under the path composition. We call this group \dfntxt{the fundamental group} of $X$ at $x_0$.
\end{definition}

\begin{theorem}{}{}
Let $\tau$ be a path in a space $X$ with $a,b$ as its initial and terminal points, respectively. Then the assignment $[\omega] \mapsto [\tau^{-1} * \omega * \tau]$ defines an isomorphism $h_{[\tau]}: \pi_1(X, a) \rightarrow \pi_1(X,b)$.
\end{theorem}

\begin{definition}{}{}
    Given a map $f:(X,a)\longrightarrow(Y,b)$, we have a homomorphism $f_{\#}:\pi_{1}(X,a)\longrightarrow\pi_{1}(Y,b)$ defined by
\begin{equation}
f_{\sharp }[\omega]=[f\circ\omega].
\end{equation}
This is well defined for if $\omega\simeq\tau$ then $f\circ\omega\simeq f\circ\tau$. Since $f\circ(\omega*\tau)=(f\circ\omega)*(f\circ\tau)$, it follows that $f_{\#}$ is a homomorphism. Moreover, we have the following two properties which are verified directly.
\begin{enumerate}
\item For any space $X$, if $Id:X\longrightarrow X$ denotes the identity map then the induced homomorphism on the fundamental groups is the identity homomorphism, i.e., $Id_{\#}=Id$.
\item Given maps $f:(X,a)\longrightarrow(Y,b)$ and $g:(Y,b)\longrightarrow(Z,c)$ we have $(g\circ f)_{\#}=g_{\#}\circ f_{\#}$.
\end{enumerate}
\end{definition}
\end{document}